\documentclass[UTF8,12pt]{ctexart}
\usepackage[a4paper,margin=24mm]{geometry}
\usepackage{amsmath,amssymb,bm,graphicx,adjustbox}
\newcommand{\fitmath}[1]{\begin{adjustbox}{max width=\linewidth}$\displaystyle #1$\end{adjustbox}}
\setlength{\parindent}{0pt}
\setlength{\parskip}{.7em}
\sloppy
\allowdisplaybreaks
\begin{document}
\section*{2009 年考研数学三 · 参考答案}
\subsection*{原 PDF 第 55 页}

\subsection*{2009年真题参考答案}

\subsection*{一、选择题}

（1）C。 （2）A。 （3）A。 （4）D。 （5）B。 （6）A。 （7）D。 （8）B。


\subsection*{二、填空题}

（9）\(\displaystyle \dfrac{\displaystyle 3}{\displaystyle 2}e\)。 （10）\(\displaystyle 2\ln2+1\)。 （11）\(\displaystyle \dfrac{\displaystyle 1}{\displaystyle e}\)。 （12）8000。 （13）2。 （14）\(\displaystyle np^2\)。


\subsection*{三、解答题}

（15）\(\displaystyle f(x,\allowbreak y)\) 在点 \(\displaystyle \left(0,\allowbreak \dfrac{\displaystyle 1}{\displaystyle e}\right)\) 处取得极小值 \(\displaystyle f\left(0,\allowbreak \dfrac{\displaystyle 1}{\displaystyle e}\right)=-\dfrac{\displaystyle 1}{\displaystyle e}\)。


（16）\(\displaystyle x\ln\left(1+\sqrt{\dfrac{\displaystyle 1+x}{\displaystyle x}}\right)+\dfrac{\displaystyle 1}{\displaystyle 2}\ln(\sqrt{1+x}+\sqrt x)-\dfrac{\displaystyle \sqrt x}{\displaystyle 2(\sqrt{1+x}+\sqrt x)}-C\)，其中 \(\displaystyle C\) 为任意常数。


（17）\(\displaystyle -\dfrac{\displaystyle 8}{\displaystyle 3}\)。


（18）证明略。


（19）\(\displaystyle x=\dfrac{\displaystyle 1}{\displaystyle 3\sqrt y}+\dfrac{\displaystyle 2}{\displaystyle 3}y,\allowbreak \ x\ge1\)。


（20）（I）满足 \(\displaystyle A\boldsymbol\xi_2=\boldsymbol\xi_1\) 的所有向量为 \(\displaystyle \boldsymbol\xi_2=k_1(1,\allowbreak -1,\allowbreak 2)^{\mathrm T}+(0,\allowbreak 0,\allowbreak 1)^{\mathrm T}\)，其中 \(\displaystyle k_1\) 为任意常数；满足 \(\displaystyle A^2\boldsymbol\xi_3=\boldsymbol\xi_1\) 的所有向量为 \(\displaystyle \boldsymbol\xi_3=k_2(1,\allowbreak -1,\allowbreak 0)^{\mathrm T}+k_3(0,\allowbreak 0,\allowbreak 1)^{\mathrm T}+\left(-\dfrac{\displaystyle 1}{\displaystyle 2},\allowbreak 0,\allowbreak 0\right)^{\mathrm T}\)，其中 \(\displaystyle k_2,\allowbreak k_3\) 为任意常数。
（II）证明略。


（21）（I）\(\displaystyle a,\allowbreak a-2,\allowbreak a+1\)。
（II）\(\displaystyle a=2\)。


（22）（I）当 \(\displaystyle f_X(x)\ne0\)，即 \(\displaystyle x>0\) 时，\(\displaystyle f_{Y\mid X}(y\mid x)=\dfrac{\displaystyle f(x,\allowbreak y)}{\displaystyle f_X(x)}=\)


\[\fitmath{\displaystyle \begin{cases}\dfrac{\displaystyle 1}{\displaystyle x},&0<y<x,\\0,&\text{其他}.\end{cases}}\]


（II）\(\displaystyle \dfrac{\displaystyle e-2}{\displaystyle e-1}\)。


（23）（I）\(\displaystyle \dfrac{\displaystyle 4}{\displaystyle 9}\)。
（II）\(\displaystyle X\) 和 \(\displaystyle Y\) 的联合分布律为


\[\fitmath{\displaystyle \begin{array}{c|ccc}Y\backslash X&0&1&2\\\hline0&\dfrac{\displaystyle 1}{\displaystyle 4}&\dfrac{\displaystyle 1}{\displaystyle 6}&\dfrac{\displaystyle 1}{\displaystyle 36}\\1&\dfrac{\displaystyle 1}{\displaystyle 3}&\dfrac{\displaystyle 1}{\displaystyle 9}&0\\2&\dfrac{\displaystyle 1}{\displaystyle 9}&0&0\end{array}}\]


\clearpage
\end{document}
